\documentclass{article}
\usepackage[margin=1in, a4paper]{geometry}
\usepackage{amsmath, amssymb, amsfonts}
\usepackage{graphicx}
\usepackage{xcolor}
\usepackage{listings}
\usepackage{booktabs}
\usepackage[utf8]{inputenc}
\usepackage[T1]{fontenc}
\usepackage{hyperref}
\hypersetup{colorlinks=true, urlcolor=blue, linkcolor=blue}
\usepackage{enumitem}
\usepackage{float}

\definecolor{codegreen}{rgb}{0,0.6,0}
\definecolor{codegray}{rgb}{0.5,0.5,0.5}
\definecolor{codepurple}{rgb}{0.58,0,0.82}
\definecolor{backcolour}{rgb}{0.99,0.99,0.99}
% Professional color palette for academic publishing
\definecolor{myblue}{cmyk}{0.8,0.4,0,0.1}
\definecolor{mygreen}{cmyk}{0.7,0,0.5,0.2}
\definecolor{myorange}{cmyk}{0,0.5,0.8,0.1}
\definecolor{mygray}{cmyk}{0,0,0,0.4}
\definecolor{arrowcolor}{cmyk}{0.9,0.5,0,0.3}
\definecolor{nodecolor}{cmyk}{0.6,0.2,0,0.05}
\definecolor{framecolor}{cmyk}{0.4,0.1,0,0.1}

\lstdefinestyle{mystyle}{
    backgroundcolor=\color{backcolour},   
    commentstyle=\color{codegreen},
    keywordstyle=\color{magenta},
    numberstyle=\tiny\color{codegray},
    stringstyle=\color{codepurple},
    basicstyle=\ttfamily\footnotesize,
    breakatwhitespace=false,         
    breaklines=true,                 
    captionpos=b,                    
    keepspaces=true,                 
    numbers=left,                    
    numbersep=5pt,                  
    showspaces=false,                
    showstringspaces=false,
    showtabs=false,                  
    tabsize=2
}
\lstset{style=mystyle}

\title{\textbf{The Logistics \& Supply Chain Problem-Solving Playbook}\\
\large Chapter 3: The Manual Reefer Monitoring Problem}
\author{The Port \& Container Terminal Playbook}
\date{}

\begin{document}

\maketitle

\section{The Manual Reefer Monitoring Problem}

The manual inspection of refrigerated containers represents one of the most critical yet challenging operational tasks in modern container terminals. With the global cold chain logistics market exceeding \$340 billion annually, ensuring the integrity of temperature-sensitive cargo through systematic monitoring has become paramount. However, the traditional manual approach to reefer monitoring presents significant operational inefficiencies, safety risks, and compliance challenges that demand sophisticated optimization solutions.

\subsection{The Scenario: Crisis in the Cold Chain}

Terminal supervisor Maria Rodriguez faces a daunting challenge at the Port of Valencia's Container Terminal. With 2,400 refrigerated containers currently in the yard, her team of 18 inspectors must perform mandatory temperature checks every 4 hours around the clock. The recent implementation of stricter EU regulations requires detailed documentation of every inspection, including timestamp, container ID, temperature readings, alarm status, and inspector identification.

The current manual process is plagued with inefficiencies. Inspectors spend an average of 3.2 minutes per container, including travel time between units, but this varies dramatically based on container accessibility. Containers stacked three high require specialized lifting equipment, while those in remote yard sections involve significant walking distances. Weather conditions further complicate operations---during recent Mediterranean storms, inspection times doubled due to safety protocols, causing dangerous delays in temperature monitoring.

The consequences of inspection failures are severe. Last month, a pharmaceutical shipment worth \$2.8 million was lost when a malfunctioning reefer went undetected for 12 hours. The manual inspection system's inherent delays and human error potential threaten both cargo integrity and terminal liability. Maria needs an optimization framework that can schedule inspections efficiently while minimizing safety risks and ensuring regulatory compliance.

\subsection{Tier 1: The Pen \& Paper Method (Mathematical Formulation)}

The manual reefer monitoring problem can be formulated as a multi-objective vehicle routing problem with time windows and priority constraints. This mathematical framework captures the essential trade-offs between inspection efficiency, safety considerations, and cargo protection requirements.

\textbf{Sets and Parameters:}

Let $R = \{1, 2, \ldots, n\}$ represent the set of reefer containers requiring inspection, where $n$ is the total number of containers in the terminal. Each container $i \in R$ has an associated priority weight $w_i$ based on cargo value and temperature sensitivity. High-value pharmaceutical containers receive $w_i = 10$, while standard food products receive $w_i = 5$, and less critical cargo receives $w_i = 1$.

The set of inspectors is denoted as $I = \{1, 2, \ldots, m\}$, where each inspector $j \in I$ has a maximum working capacity of $C_j$ containers per shift. The inspection time for container $i$ by inspector $j$ is $t_{ij}$, which varies based on container accessibility and inspector experience level.

\textbf{Decision Variables:}

The primary decision variable is $x_{ij} \in \{0, 1\}$, where $x_{ij} = 1$ if inspector $j$ is assigned to inspect container $i$, and 0 otherwise. The variable $y_{ijk} \in \{0, 1\}$ indicates the sequence, where $y_{ijk} = 1$ if inspector $j$ visits container $i$ immediately before container $k$.

The continuous variable $s_{ij} \geq 0$ represents the start time of inspection for container $i$ by inspector $j$, while $d_{ij} \geq 0$ captures any delay beyond the required inspection interval.

\textbf{Objective Function:}

The mathematical model seeks to minimize a weighted combination of total inspection time, safety risks, and cargo protection penalties:

\begin{align}
\text{Minimize: } Z &= \alpha \sum_{i \in R} \sum_{j \in I} t_{ij} x_{ij} + \beta \sum_{i \in R} \sum_{j \in I} r_{ij} x_{ij} + \gamma \sum_{i \in R} w_i d_{ij}
\end{align}

where $r_{ij}$ represents the safety risk score for inspector $j$ accessing container $i$, and $\alpha$, $\beta$, $\gamma$ are weighting parameters balancing efficiency, safety, and cargo protection objectives.

\textbf{Constraints:}

Every container must be assigned to exactly one inspector:
\begin{align}
\sum_{j \in I} x_{ij} &= 1, \quad \forall i \in R
\end{align}

Inspector capacity limitations:
\begin{align}
\sum_{i \in R} t_{ij} x_{ij} &\leq C_j, \quad \forall j \in I
\end{align}

Time window constraints for mandatory inspection intervals:
\begin{align}
s_{ij} + d_{ij} &\leq L_i x_{ij}, \quad \forall i \in R, j \in I
\end{align}

where $L_i$ is the maximum allowable time between inspections for container $i$.

\textbf{Visualization of the Problem Structure:}

\begin{figure}[htbp]
\centering
\includegraphics[width=\textwidth]{../figures-png/line3.fig1.png}
\caption{Mathematical Model Visualization: Container Terminal with Inspector Routes and Safety Zones}
\end{figure}

\textbf{Concrete Example Solution:}

Consider a simplified scenario with 6 containers and 2 inspectors working a 4-hour shift. Container priorities and locations are:
- Container 1: Pharmaceutical (Priority 10), Ground level, $t_{11} = 2$ min, $t_{12} = 3$ min
- Container 2: Food products (Priority 5), Stack level 2, $t_{21} = 5$ min, $t_{22} = 4$ min  
- Container 3: General cargo (Priority 1), Ground level, $t_{31} = 2$ min, $t_{32} = 2$ min
- Container 4: Pharmaceutical (Priority 10), Stack level 3, $t_{41} = 8$ min, $t_{42} = 6$ min
- Container 5: Food products (Priority 5), Ground level, $t_{51} = 2$ min, $t_{52} = 3$ min
- Container 6: General cargo (Priority 1), Stack level 2, $t_{61} = 5$ min, $t_{62} = 4$ min

Inspector capacities: $C_1 = 120$ minutes, $C_2 = 120$ minutes.

The optimal solution assigns:
- Inspector 1: Containers 1, 3, 5 (total time: 6 minutes, all ground level for safety)
- Inspector 2: Containers 2, 4, 6 (total time: 14 minutes, experienced with stacked containers)

This assignment minimizes total weighted completion time while respecting safety protocols by assigning the most dangerous high-stack inspections to the more experienced inspector.

\subsection{Tier 2: The Classic Heuristic (Python Implementation)}

The manual reefer monitoring problem requires a practical heuristic approach that can quickly generate high-quality inspection schedules in real-time terminal operations. We implement a priority-based greedy algorithm that sequentially assigns containers to inspectors based on a composite scoring function.

\textbf{Algorithm Design Philosophy:}

The greedy heuristic operates on the principle of immediate optimization, making locally optimal choices at each step. The algorithm evaluates each unassigned container using a multi-criteria scoring function that balances cargo priority, inspection urgency, safety considerations, and inspector workload balance.

\textbf{Detailed Pseudocode:}

\lstinputlisting[language=Python, caption=Priority-Based Greedy Reefer Inspection Algorithm]{.././codes-python/line3.py1.py}

\textbf{Algorithm Visualization:}

\begin{figure}[htbp]
\centering
\includegraphics[width=\textwidth]{../figures-png/line3.fig2.png}
\caption{Enhanced Priority-Based Greedy Algorithm with Node-Based Flow Structure and Performance Metrics}
\end{figure}

\textbf{Time Complexity Analysis:}

The greedy algorithm has a time complexity of $O(n^2 \cdot m)$ where $n$ is the number of containers and $m$ is the number of inspectors. For each of the $n$ containers to be assigned, the algorithm evaluates $n \cdot m$ possible assignments, resulting in the quadratic behavior. The subsequent 2-opt optimization adds $O(k^3)$ complexity per inspector, where $k$ is the average number of containers per inspector, but this is typically much smaller than the main assignment loop.

\textbf{Concrete Example Output:}

Running the algorithm on a terminal with 24 reefer containers and 3 inspectors produces the following optimized schedule:

\lstinputlisting[language=Python, caption=Sample Algorithm Output]{.././codes-python/line3.py2.py}

The algorithm successfully balances multiple objectives, assigning high-priority pharmaceutical containers to be inspected first while distributing dangerous high-stack inspections to the most experienced inspector. The 2-opt optimization reduces total travel time by an average of 12\%, significantly improving operational efficiency.

\subsection{Tier 3: The Advanced Algorithm (Metaheuristic Implementation)}

The complexity of real-world reefer monitoring demands a sophisticated metaheuristic approach capable of escaping local optima while handling multiple conflicting objectives. We implement a Simulated Annealing algorithm enhanced with adaptive cooling schedules and multi-neighborhood search strategies to optimize inspection schedules under dynamic terminal conditions.

\textbf{Simulated Annealing Theory:}

Simulated Annealing draws inspiration from the metallurgical process where controlled cooling allows metal atoms to reach low-energy crystalline structures. In optimization terms, the algorithm accepts worse solutions with a probability that decreases over time, enabling exploration of the solution space while gradually converging to high-quality solutions.

The acceptance probability follows the Metropolis criterion:
\begin{align}
P(\text{accept}) = \begin{cases}
1 & \text{if } \Delta E \leq 0 \\
e^{-\Delta E / T} & \text{if } \Delta E > 0
\end{cases}
\end{align}

where $\Delta E$ represents the change in objective function value and $T$ is the current temperature parameter.

\textbf{Enhanced Algorithm Design:}

\lstinputlisting[language=Python, caption=Adaptive Simulated Annealing for Reefer Monitoring]{.././codes-python/line3.py3.py}

\textbf{Metaheuristic Visualization:}

\begin{figure}[htbp]
\centering
\includegraphics[width=\textwidth]{../figures-png/line3.fig3.png}
\caption{Simulated Annealing Temperature Schedule and Solution Convergence}
\end{figure}

\textbf{Concrete Example Performance:}

Testing the enhanced Simulated Annealing algorithm on a complex scenario with 48 reefer containers and 6 inspectors yields impressive results:

\lstinputlisting[language=Python, caption=Metaheuristic Algorithm Results]{.././codes-python/line3.py4.py}

The Simulated Annealing algorithm demonstrates superior performance compared to the greedy approach, achieving a 28.5\% improvement in solution quality while maintaining perfect compliance with safety and timing constraints. The adaptive mechanisms successfully balance exploration and exploitation, resulting in robust solutions that handle the complex multi-objective nature of the reefer monitoring problem.

\subsection{Tier 4: The AI/ML/RL Augmentation Method}

The dynamic nature of container terminal operations and the continuous influx of real-time data make Reinforcement Learning an ideal paradigm for automated reefer inspection scheduling. We develop a Deep Q-Network (DQN) agent that learns optimal inspection policies through interaction with a simulated terminal environment, adapting to changing conditions and improving performance over time.

\textbf{Reinforcement Learning Problem Formulation:}

The reefer monitoring problem maps naturally to a Markov Decision Process (MDP) where an intelligent agent must make sequential decisions about inspection assignments while learning from the consequences of its actions.

\textbf{State Space Definition ($S$):}

The state space captures the complete terminal status at any decision point:
\begin{align}
s_t = \{C_t, I_t, E_t, T_t, W_t\}
\end{align}

where:
- $C_t$: Container status vector (temperature, time since last inspection, priority, location)
- $I_t$: Inspector availability and location vector
- $E_t$: Environmental conditions (weather, equipment status, traffic)
- $T_t$: Current time and scheduling constraints
- $W_t$: Workload distribution and fatigue levels

\textbf{Action Space Definition ($A$):}

Actions represent inspection assignment decisions:
\begin{align}
a_t = (\text{container\_id}, \text{inspector\_id}, \text{inspection\_time})
\end{align}

The action space includes all valid assignments plus special actions for rescheduling and emergency responses.

\textbf{Reward Function Design ($R$):}

The reward function balances multiple objectives with carefully designed components:
\begin{align}
R(s_t, a_t, s_{t+1}) = R_{\text{efficiency}} + R_{\text{safety}} + R_{\text{compliance}} + R_{\text{penalty}}
\end{align}

where:
- $R_{\text{efficiency}} = -\alpha \cdot \text{inspection\_time} - \beta \cdot \text{travel\_time}$
- $R_{\text{safety}} = \gamma \cdot \text{safety\_score\_improvement}$
- $R_{\text{compliance}} = \delta \cdot \text{compliance\_bonus} - \epsilon \cdot \text{violation\_penalty}$
- $R_{\text{penalty}} = -\zeta \cdot \text{workload\_imbalance}$

\textbf{Deep Q-Network Implementation:}

\lstinputlisting[language=Python, caption=DQN Agent for Reefer Inspection Scheduling]{.././codes-python/line3.py5.py}

\textbf{RL Algorithm Visualization:}

\begin{figure}[htbp]
\centering
\includegraphics[width=\textwidth]{../figures-png/line3.fig4.png}
\caption{Deep Q-Network Architecture and Learning Components}
\end{figure}

\textbf{Concrete Training Results:}

\lstinputlisting[language=Python, caption=DQN Training Performance Results]{.././codes-python/line3.py6.py}

The Deep Q-Network agent demonstrates remarkable adaptation capabilities, learning complex inspection policies that significantly outperform traditional scheduling methods. The agent discovers non-obvious strategies such as strategic delays during adverse weather and proactive workload balancing that human schedulers often miss. After deployment, the system continues learning and improving, showing the dynamic optimization capabilities that make reinforcement learning particularly suitable for the complex, evolving environment of container terminal operations.

\section{Practice Questions}

\textbf{Tier 1: Mathematical Formulation Problems}

1. \textbf{Basic Integer Programming Model:} A container terminal has 12 reefer containers and 3 inspectors available for the next shift. Container priorities are: 4 pharmaceutical (weight 10), 5 food products (weight 5), and 3 general cargo (weight 1). Inspection times vary by container accessibility: ground level (2 min), stack level 2 (4 min), stack level 3 (7 min). Each inspector has a 4-hour (240 min) capacity. Formulate the integer programming model to minimize total weighted inspection time while ensuring all containers are inspected within their priority windows (pharmaceutical: 2 hours, food: 3 hours, general: 4 hours).

2. \textbf{Multi-Objective Optimization:} Extend the basic model to include safety constraints where stack level 3 containers can only be inspected by certified inspectors (2 out of 3 available). Add workload balance constraints ensuring no inspector is assigned more than 150\% of the average workload. Reformulate as a goal programming problem with priorities: (1) Safety compliance, (2) Time window adherence, (3) Workload balance, (4) Efficiency.

3. \textbf{Dynamic Programming Formulation:} Consider a sequential inspection problem where 6 containers must be inspected in order, but inspectors can be switched between containers. Each container-inspector pair has different costs based on travel time and expertise match. Formulate the dynamic programming recursion and solve for the optimal inspector assignment sequence. Given data: Container sequence A→B→C→D→E→F with inspector costs matrix where element (i,j) represents cost for inspector i to handle container j.

\textbf{Tier 2: Heuristic Algorithm Problems}

4. \textbf{Greedy Algorithm Implementation:} Implement and test a modified greedy algorithm for 18 containers (6 high-priority, 7 medium-priority, 5 low-priority) and 4 inspectors with different skill levels. The algorithm should use a composite scoring function combining urgency (40\%), safety (35\%), and workload balance (25\%). Provide the complete assignment solution and calculate total completion time. Include 2-opt local improvement results.

5. \textbf{Constructive Heuristic Design:} Design a priority-based constructive heuristic that builds inspection routes considering spatial clustering of containers. Given a terminal layout with containers distributed across 8 zones, each zone containing 3-6 reefers, develop a zone-based assignment strategy. Calculate and compare the solution quality against a simple nearest-neighbor approach using provided distance matrix data.

6. \textbf{Multi-Start Improvement Heuristic:} Develop a multi-start local search algorithm that generates 10 different initial solutions and applies variable neighborhood descent. Test on a scenario with 30 containers, 5 inspectors, and compare final solution quality. Include analysis of solution diversity and convergence characteristics.

\textbf{Tier 3: Metaheuristic Algorithm Problems}

7. \textbf{Simulated Annealing Parameter Tuning:} Design and tune a Simulated Annealing algorithm for a large-scale problem with 60 containers and 8 inspectors. Test different cooling schedules (linear, geometric, logarithmic) and neighborhood operators (swap, relocate, reverse). Provide parameter sensitivity analysis and recommend optimal configuration. Expected improvement over greedy heuristic should be 15-25\%.

8. \textbf{Genetic Algorithm Design:} Implement a genetic algorithm with specialized crossover and mutation operators for the inspection scheduling problem. Design chromosome representation for variable-length inspector routes, develop order-preserving crossover, and implement adaptive mutation rates. Test on 45 containers with 6 inspectors and analyze population diversity metrics over 200 generations.

9. \textbf{Hybrid Metaheuristic Development:} Create a hybrid algorithm combining Particle Swarm Optimization with Variable Neighborhood Search. The PSO component should explore different workload distributions while VNS refines individual inspector routes. Demonstrate on a 72-container, 9-inspector scenario and show convergence characteristics. Compare with pure PSO and pure VNS implementations.

\textbf{Tier 4: AI/ML/RL Implementation Problems}

10. \textbf{Deep Reinforcement Learning Environment:} Implement a custom OpenAI Gym environment for reefer inspection scheduling with continuous state space (container temperatures, inspector locations, weather conditions) and discrete action space. Design appropriate reward shaping to encourage both efficiency and safety. Train a DQN agent for 1000 episodes and analyze learned policies. Expected final performance should achieve 90\%+ of optimal benchmark solution.

11. \textbf{Supervised Learning Prediction Model:} Develop a machine learning model to predict inspection duration based on historical data. Features include container type, stack level, inspector experience, weather conditions, and time of day. Use provided dataset of 5000+ inspection records to train Random Forest and Neural Network models. Evaluate prediction accuracy and integrate into scheduling optimization framework.

12. \textbf{Multi-Agent Reinforcement Learning:} Design a multi-agent RL system where each inspector is an autonomous agent learning optimal inspection strategies. Implement cooperation mechanisms for workload sharing and conflict resolution. Test coordination effectiveness in scenarios with 8 agents and 96 containers. Analyze emergent behaviors and compare with centralized scheduling approaches. Demonstrate adaptation to dynamic priority changes and equipment failures.

\end{document}
